% TOP_PROBLEM: {"id": 4300004, "title": "Typical group automorphisms", "category_id": 1, "category": {"id": 1, "name": "number_theory", "display_name": "Number Theory", "description": "Properties of integers, prime numbers, Diophantine equations.", "slug": "number-theory", "order_index": 1, "created_at": "2026-07-31T15:26:25.670Z"}, "status": "partially_solved", "problem_number": "AMR-042-0004", "set_id": 13, "statusQualification": "The logarithm encloses the full periodic-point product; this resolves the missing-parentheses ambiguity in the catalogue transcription."}
% FIRST_POSTED: 2026-10-02
% ENTRY_KIND: statement_only
% UnsolvedMath ID 4300004; source catalogue code AMR-042-0004
% !TeX program = lualatex
% Definition and sources only; this file is not an LLM solution attempt.
\documentclass[11pt,a4paper]{article}
\usepackage[margin=25mm]{geometry}
\usepackage{fontspec}
\setmainfont{lmroman10-regular.otf}[BoldFont=lmroman10-bold.otf,
 ItalicFont=lmroman10-italic.otf,BoldItalicFont=lmroman10-bolditalic.otf]
\usepackage{amsmath,amssymb,mathtools}
\usepackage{unicode-math}
\setmathfont{latinmodern-math.otf}
\AtBeginDocument{\let\setminus\smallsetminus}
\usepackage{xurl,hyperref}
\hypersetup{colorlinks=true,urlcolor=blue}
\setcounter{secnumdepth}{0}
\setlength{\parindent}{0pt}
\setlength{\parskip}{5pt}
\setlength{\emergencystretch}{3em}
\begin{document}
\section*{Typical group automorphisms}\label{4300004}
{\small UnsolvedMath ID 4300004 · AMR-042-0004 · Number Theory · AMR Open Problem Lists}
\subsection{Definitions and mathematical statement}
For each prime $p$, independently choose a Bernoulli random variable
$\xi_p\in\{0,1\}$ with probability $1/2$ of either value, and set
$Q=\{p:\xi_p=1\}$. For a positive integer $a$, let $v_p(a)$ be the
largest integer $j\ge0$ such that $p^j$ divides $a$, and put
$|a|_p=p^{-v_p(a)}$. Define
\[
 F_Q(n)=(2^n-1)\prod_{p\in Q}|2^n-1|_p
       =\prod_{p\notin Q}p^{v_p(2^n-1)}.
\]
The products are finite in effect, because only prime divisors of
$2^n-1$ contribute a nonunit factor. The intended periodic-point growth
question is whether
\[
 \mathbb P\!\left(\limsup_{n\to\infty}\frac{\log F_Q(n)}{n}
                      =\log2\right)=1.
\]
The same random set $Q$ is used for every $n$. All logarithms are real
natural logarithms. Since $1\le F_Q(n)\le2^n-1$, the claimed value is the
largest possible exponential upper growth rate.

Parentheses in the formula matter: the logarithm applies to the entire
product defining $F_Q(n)$, as required by the periodic-point interpretation
and the source's lower bound. This is stronger than computing an
expectation for each $n$: almost every single infinite prime selection
must have a subsequence along which its retained prime powers recover
the full exponential growth rate. No independence between the numbers
$F_Q(n)$ for different $n$ is assumed. In particular, repeating the prime
selection afresh for each index would define a different probability problem.
\subsection{Short English statement}
For almost every independent half-selection of primes, do the surviving factors of $2^n-1$ retain the full exponential limsup growth?
\subsection{Sources}
\begin{itemize}
\item[S1] ULAM.AI and the UnsolvedMath Contributors, supplied problems.json, record 4300004 (AMR-042-0004), AMR Open Problem Lists. Catalogue identity and source transcription; CC BY 4.0.
\url{https://huggingface.co/datasets/ulamai/UnsolvedMath}.
\item[S2] Ward - Six problems in Algebraic Dynamics (2006), Problem D. Original source indexed by AMR-042-0004.
\url{https://www.imath.kiev.ua/~skolyada/kevin.pdf}.
\end{itemize}
\end{document}
